\chapter[Metodologia]{Metodologia}
\label{cap:Metodologia}

A metodologia adotada cobre cinco modelos de comparação principal e duas linhas de base multimodais usadas como estudo exploratório complementar. Os cinco modelos principais atuam sobre a mesma tarefa, a verificação binária de parentesco no benchmark Families in the Wild, derivado do Track-I oficial do RFIW e avaliado por validação cruzada de cinco dobras em nível de família. O conjunto de comparação inclui a linha de base sem treinamento com extrator facial pronto (B0, AdaFace IR-101 com cosseno), um ViT com atenção cruzada bidirecional FaCoR (Modelo 02), um DINOv2 com adaptação por LoRA e atenção cruzada diferencial (Modelo 05), um modelo aumentado por recuperação sobre galeria de pares positivos (Modelo 06) e a proposta autoral do Modelo 12, com AdaFace parcialmente descongelado, tokens regionais anatômicos, cabeça auxiliar de relação e passagem simétrica. O Modelo 02 também é avaliado em KinFaceW-I com validação cruzada de cinco dobras, como ponto de comparação adicional com a literatura clássica do banco menor. Dois modelos periféricos do projeto, o Modelo 01 (linha de base histórica com CNN e síntese de idade) e o Modelo 03 (estudo de capacidade do híbrido ConvNeXt mais ViT), ficam descritos no Apêndice~\ref{ap:modelos-perifericos}. As duas linhas de base com modelos multimodais generalistas tratam uma tarefa diferente, classificação multiclasse das onze relações do FIW, e por isso aparecem em seção própria como avaliação exploratória, não como competidoras diretas dos modelos supervisionados.

\section{Desenho da pesquisa}

O problema de verificação de parentesco facial é tratado neste trabalho como pesquisa experimental e comparativa. Cinco modelos supervisionados são treinados e avaliados sob o mesmo protocolo, e duas linhas de base com modelos multimodais generalistas são executadas sobre o mesmo conjunto de teste, em regime zero-shot e em adaptação por prompt. A comparação principal é feita por ROC AUC, métrica independente de limiar, complementada por Average Precision, TAR@FAR=0,1, 0,01 e 0,001, e acurácia por tipo de relação. A análise por tipo de relação acompanha as métricas globais, porque a literatura mostra que relações de segundo grau entre avós e netos comportam-se de forma muito diferente das relações de primeiro grau entre pais e filhos ou entre irmãos.

Três dos cinco modelos servem como referência e cobrem famílias de método consolidadas. Os outros dois são exploratórios e combinam técnicas recentes ainda pouco validadas para parentesco facial na literatura consultada. As linhas de base multimodais entram apenas como medida do que um VLM genérico consegue inferir, e não competem diretamente com os modelos treinados, porque avaliam tarefa diferente (classificação multiclasse das 11 relações, não verificação binária de parentesco).

\section{Datasets}

\subsection{KinFaceW-I}

KinFaceW-I é um dos bancos clássicos para verificação de parentesco facial. O conjunto tem cerca de 1.066 pares de imagens distribuídos em quatro tipos de relação, sendo pai-filha, pai-filho, mãe-filha e mãe-filho. As imagens são frontais, com fundo controlado, e foram coletadas de fontes públicas com curadoria manual. Neste trabalho, KinFaceW-I serve como banco de prototipagem rápida e como base da validação cruzada de cinco dobras, com splits disjuntos por par.

\subsection{Families in the Wild}

Families in the Wild, ou FIW, é hoje o maior banco público de parentesco facial. A versão usada neste trabalho é a do RFIW Track-I, com 761 famílias e onze relações cobertas. As relações incluem pai-filho e pai-filha (fs, fd), mãe-filho e mãe-filha (ms, md), irmão-irmão e irmã-irmã (bb, ss), irmãos de gêneros distintos (sibs), e quatro relações de avós e netos (gfgs, gfgd, gmgs, gmgd). As imagens cobrem uma variedade muito maior de condições do que KinFaceW-I, com diferenças de iluminação, pose, idade e qualidade. O FIW é o benchmark principal deste trabalho e a avaliação é conduzida por validação cruzada de cinco dobras em nível de família, preservando a premissa do Track-I de que indivíduos de uma mesma família não podem aparecer simultaneamente em treino, validação e teste. Em cada dobra, as famílias são separadas em treino, validação e teste, e os pares positivos e negativos variam conforme a composição familiar de cada partição.

\subsection{Amostragem de pares e negativos}

O número de pares positivos disponíveis nos dados é finito e pequeno em comparação com o número possível de pares negativos. O carregador de dados amostra pares negativos restringindo o sorteio a indivíduos de famílias distintas, garantia herdada da divisão disjunta por família do Track-I do RFIW. Essa restrição reduz o risco de que pares negativos contenham parentes não anotados, porque membros de uma mesma família nunca aparecem simultaneamente em treino, validação e teste dentro da mesma dobra.

Na estratégia \texttt{random}, sorteiam-se aleatoriamente dois indivíduos de famílias diferentes, sem condicionar a amostragem ao tipo de relação. O identificador \texttt{relation\_matched} aparece no log experimental como uma tentativa de preservar, no par negativo, a classe de relação, o gênero e uma faixa etária aproximada do par positivo de referência. Uma auditoria do ciclo R004 do Modelo 12 mostrou que a implementação usada no FIW não preservava efetivamente relação ou papel; na prática, ela se comportava como uma nova amostragem aleatória com semente diferente. Por isso, os resultados que mencionam \texttt{relation\_matched} são interpretados como variação de semente do amostrador, não como teste válido de negativos difíceis.

A razão entre negativos e positivos é 2:1 no treino, com reamostragem por época segundo a semente fixada, e 1:1 na validação e no teste, com listas de negativos fixas para cada dobra. A semente principal da amostragem é fixada em 42, de modo que execuções comparáveis dentro de cada família de modelo possam ser reexecutadas a partir do log experimental. Quando uma execução altera o amostrador ou seu deslocamento de semente, essa diferença é tratada explicitamente na interpretação dos resultados.

\subsection{Pré-processamento e alinhamento facial}

Todas as imagens usadas pelos modelos supervisionados passam por alinhamento facial antes do treinamento. A versão principal do FIW alinhado usa detecção de faces por InsightFace \texttt{buffalo\_l} (RetinaFace) e estimação afim por \texttt{cv2.estimateAffinePartial2D} com LMEDS para projetar cinco pontos faciais no template ArcFace ampliado para 224x224 pixels. Essa resolução coincide com a entrada nativa dos backbones ViT-B/16 e DINOv2 usados nos Modelos 02, 05 e 06, evitando redimensionamentos adicionais no carregador.

Uma segunda versão alinhada, usada nos modelos derivados da linhagem ArcFace/AdaFace quando necessário, emprega MTCNN da biblioteca \texttt{facenet-pytorch} e transformação de similaridade de Umeyama para o mesmo template ArcFace em 224x224 pixels. Quando o detector falha, o procedimento aplica recorte central quadrado e redimensionamento, preservando a validade das listas de pares e evitando descarte silencioso de amostras. O Modelo 12 recebe a imagem alinhada em 224x224 e, a partir dela, recorta cinco regiões anatômicas fixas que são redimensionadas para 112x112 antes da extração de embeddings regionais pelo AdaFace.

Para os VLMs, cada entrada visual é composta por duas faces já alinhadas, posicionadas lado a lado em uma única imagem. A composição preserva a ordem do par registrada no FIW para que a relação multiclasse (por exemplo, pai-filho ou mãe-filha) permaneça semanticamente definida.

\begin{table}[h]
\centering
\caption{Resumo dos datasets utilizados.}
\label{tab:datasets}
\begin{tabular}{lrrrl}
\hline
\textbf{Banco} & \textbf{Famílias} & \textbf{Pares} & \textbf{Relações} & \textbf{Papel no trabalho} \\
\hline
KinFaceW-I & --- & {\raise.17ex\hbox{$\scriptstyle\sim$}}1.066 & 4 & Prototipagem e validação cruzada de cinco dobras \\
Families in the Wild & 761 & variável por dobra\footnotemark & 11 & Benchmark principal em validação cruzada por família \\
\hline
\end{tabular}
\footnotetext{Os pares positivos e negativos variam entre dobras porque a partição é feita por família e preserva a disjunção familiar do Track-I do RFIW. No treino, a proporção é 2:1 entre negativos e positivos; em validação e teste, a proporção é aproximadamente 1:1.}
\end{table}

\section{Modelos avaliados}

Os modelos descritos a seguir correspondem ao eixo central da comparação no Capítulo~\ref{cap:resultados}. O B0 entra como referência sem treinamento e linha de base mínima recomendada pelo documento de referências do projeto. Os Modelos 02, 05 e 06 são as três famílias supervisionadas comparadas, e o Modelo 12 é a proposta principal deste TCC. A Tabela~\ref{tab:modelos_metodologia} resume os cinco modelos centrais por backbone, parâmetros treináveis e ideia central. Dois modelos adicionais foram conduzidos ao longo do projeto e ficam fora do eixo de comparação principal: o Modelo 01 atua como linha de base histórica herdada da versão de 2018 do pré-projeto, e o Modelo 03 atua como estudo de capacidade da fusão híbrida ConvNeXt mais ViT. Suas descrições e resultados ficam no Apêndice~\ref{ap:modelos-perifericos}.

A numeração dos modelos preserva o histórico experimental do repositório, e por isso não é sequencial no corpo do texto. O Modelo 04 não gerou execução consolidada para o eixo comparativo. Os Modelos 07, 08, 09, 10 e 11 foram iterações intermediárias com backbones ArcFace ou AdaFace, úteis para diagnóstico interno de treinamento, mas não são tratados como competidores finais porque suas hipóteses foram absorvidas ou superadas pelo desenho do Modelo 12. Essa decisão evita misturar protótipos descartados com as arquiteturas selecionadas para a comparação principal.

\begin{table}[h]
\centering
\caption{Visão geral dos cinco modelos no eixo central de comparação, com backbone, total de parâmetros treináveis e ideia central de cada um.}
\label{tab:modelos_metodologia}
\begin{tabular}{llrl}
\hline
\textbf{Modelo} & \textbf{Backbone} & \textbf{Params treináveis} & \textbf{Ideia central} \\
\hline
B0  & AdaFace IR-101 (congelado)          & 0 & Extrator facial pronto com cosseno e limiar em validação \\
M02 & vit\_base\_patch16\_224              & {\raise.17ex\hbox{$\scriptstyle\sim$}}86M & Atenção cruzada bidirecional entre faces \\
M05 & DINOv2 ViT-B/14 + LoRA / descongelamento & 8,5M a 95M & Backbone auto-supervisionado com atenção cruzada diferencial \\
M06 & ViT-B/16 ou DINOv2 (congelado)      & {\raise.17ex\hbox{$\scriptstyle\sim$}}8-12M & Memória não-paramétrica de pares positivos \\
M12 & AdaFace IR-101 (estágio 4 + output\_layer) & 31,0M a 31,6M & Tokens regionais anatômicos com atenção cruzada entre regiões, gate e passagem simétrica \\
\hline
\end{tabular}
\end{table}

\subsection{Linha de base B0, AdaFace cosseno congelado}

O B0 implementa a linha de base mínima recomendada pelo documento de referências do projeto. O backbone é o AdaFace IR-101 \cite{kim2022adaface} pré-treinado em WebFace4M, com todos os pesos congelados. Para cada par $(x_a, x_b)$, o modelo extrai embeddings $L^2$-normalizados de 512 dimensões e calcula a similaridade cosseno. O escore é mapeado para o intervalo $[0, 1]$ por meio de $(\cos + 1)/2$, e o limiar de decisão é escolhido em validação por busca do ponto que maximiza F1. Não há treinamento e não há parâmetros treináveis. O B0 estabelece o desempenho de referência da tarefa de parentesco com um extrator facial pronto, sem qualquer adaptação, e serve para isolar o ganho arquitetônico dos modelos treinados.

\subsection{Modelo 02, ViT com atenção cruzada bidirecional}

O Modelo 02 substitui o backbone convolucional por um Vision Transformer, vit\_base\_patch16\_224, pré-treinado em ImageNet, e adiciona um módulo de atenção cruzada bidirecional inspirado no FaCoR \cite{su2023facornet}. A atenção cruzada permite que cada face atenda aos tokens de patch da outra, de modo que o modelo aprende quais regiões comparar entre as duas imagens, em vez de só comparar embeddings globais. O módulo tem duas camadas e oito cabeças, e é seguido por uma camada de channel attention no estilo squeeze-and-excitation. A saída passa por uma projeção linear que produz embeddings de 512 dimensões com normalização L2, e o modelo total tem cerca de 86 milhões de parâmetros.

\subsection{Modelo 05, DINOv2 com LoRA e atenção cruzada diferencial}

O Modelo 05 combina três técnicas recentes ainda pouco validadas para parentesco facial na literatura consultada. O backbone é o DINOv2 \cite{oquab2024dinov2} na variante vit\_base\_patch14\_dinov2, treinado de forma auto-supervisionada em 142 milhões de imagens sem rótulos. A adaptação ao domínio é feita por adaptadores LoRA \cite{hu2021lora} injetados nas projeções qkv e proj de cada bloco do Transformer, com rank 8 e alpha 16 na configuração padrão. O módulo de atenção cruzada é uma variante diferencial, proposta originalmente para modelagem de linguagem em 2024, que computa a atenção como a diferença entre duas projeções softmax e suprime correspondências espúrias entre regiões irrelevantes das faces. Uma cabeça auxiliar prediz o tipo da relação a partir de um token aprendível, em paralelo com a cabeça binária de verificação.

O Modelo 05 cobre quatro regimes de plasticidade do backbone, comparados no Capítulo 5. No regime de referência, o backbone DINOv2 permanece congelado e o treinamento atinge apenas LoRA, atenção cruzada e cabeças, com cerca de 8,5 milhões de parâmetros treináveis, aproximadamente dez vezes menos que o Modelo 02. As outras três variantes flexibilizam o backbone em graus diferentes. Quando os quatro blocos finais do DINOv2 ganham gradiente com taxa de aprendizado reduzida por um fator de cem em relação à cabeça, o descongelamento parcial chega a cerca de 30 milhões de parâmetros ativos. Liberando todo o backbone e usando pico de taxa de aprendizado no estilo do Modelo 02, o descongelamento completo ativa cerca de 95 milhões de parâmetros. Já a variante híbrida abandona o ramo único do DINOv2 e usa dois encoders congelados em paralelo, DINOv2 e o vit\_base\_patch16\_224 do Modelo 02, deixando cerca de 11,6 milhões de parâmetros treináveis apenas nos módulos de fusão e atenção. Rank LoRA, dropout, escala de taxa de aprendizado e razão entre BCE e Supervised Contrastive variam por execução e ficam detalhados no log experimental, resumidos no Capítulo 5.

\subsection{Modelo 06, aumentado por recuperação sobre galeria}

O Modelo 06 adota uma abordagem de memória não-paramétrica. Em vez de forçar a decisão de parentesco numa única passagem feed-forward, o modelo mantém uma galeria de pares positivos do conjunto de treino e, em inferência, recupera os K pares mais similares ao par de consulta. Os pares recuperados, junto com seus tipos de relação, viram tokens de contexto para um módulo de atenção cruzada que olha da consulta para os suportes. O encoder permanece congelado, e a galeria é reconstruída a cada dobra a partir da respectiva partição de treino, sem exemplos de validação ou teste. A configuração padrão usa K igual a 32 e duas camadas de atenção cruzada, com total treinável entre 8 e 12 milhões de parâmetros, dependendo da configuração de recuperação. A motivação é dar acesso explícito a exemplos conhecidos de parentesco que se parecem com o par de consulta, em vez de tentar absorver tudo nos pesos.

\subsection{Modelo 12, proposta autoral com tokens regionais e descongelamento parcial}

O Modelo 12 é a proposta principal e autoral deste trabalho. O backbone é o AdaFace IR-101 \cite{kim2022adaface} pré-treinado em WebFace4M, especializado em reconhecimento facial em larga escala. A estratégia de plasticidade do backbone mantém os três primeiros estágios congelados e libera apenas o estágio 4 (três blocos BasicBlockIR de saída) e o output\_layer, que produz o embedding de 512 dimensões. Esse descongelamento parcial adiciona cerca de 26 milhões de parâmetros treináveis sobre o esqueleto de tokens regionais e o adaptador, totalizando 31,55 milhões de parâmetros treináveis em 70,74 milhões totais (44,6\%) nas variantes completas.

Sobre o backbone facial, o Modelo 12 constrói uma representação regional explícita. Cinco regiões anatômicas são definidas em coordenadas fixas na imagem alinhada de 224x224 pixels. As regiões cobrem face inteira (global), faixa horizontal superior (olhos), região central (nariz), faixa inferior média (boca) e faixa inferior (mandíbula). Cada região é recortada e redimensionada para 112x112, e o AdaFace produz um embedding por região. Os cinco tokens passam por uma camada bidirecional de atenção cruzada entre regiões com 4 cabeças, que computa um mapa de atenção 5x5 entre as regiões das duas faces. Esse mapa permite que cada região de uma face atenda às regiões correspondentes ou complementares da outra, em escala muito menor do que os mapas de patch 196x196 dos modelos baseados em ViT.

A decisão final combina três componentes. Para cada região $k \in \{1, \ldots, K\}$ com $K = 5$, o modelo calcula a similaridade entre os tokens correspondentes das duas faces, $s_k \in \mathbb{R}$. Um módulo de gating sigmoide aprende pesos $w_k \in [0, 1]$ por região, interpretáveis como saliência regional para o par específico em análise. Diferentemente de softmax, o sigmoide permite que mais de uma região seja saliente simultaneamente. O classificador final da variante de referência recebe um vetor concatenado de dimensão $4 d + 2 K + 1 = 2059$, com $d = 512$ a dimensão do embedding do AdaFace. O vetor é formado pela concatenação de $g_A$ e $g_B$ (os tokens globais das duas faces), $|g_A - g_B|$ (diferença absoluta), $g_A \odot g_B$ (produto elemento a elemento), o vetor de similaridade regional $s \in \mathbb{R}^K$, o vetor de pesos do gate $w \in \mathbb{R}^K$, e o produto interno escalar $w^\top s$ entre os dois últimos. Em uma ablação posterior, denominada fusão apenas comparativa, os vetores globais crus $g_A$ e $g_B$ são removidos e o classificador recebe apenas primitivas de comparação, reduzindo a entrada para 1035 dimensões. A saída é um único logit, treinado por entropia cruzada binária (BCEWithLogitsLoss). A taxa de aprendizado pico é 1e-5, dez vezes menor do que a usada com backbone totalmente congelado, escolhida para proteger as representações identitárias profundas do AdaFace de serem reescritas drasticamente.

A avaliação do Modelo 12 cobre uma sequência de fases. Na primeira fase (R001), o backbone AdaFace fica totalmente congelado e apenas os 5,6 milhões de parâmetros do esqueleto regional, do adaptador, do gate e da cabeça são treinados, com taxa de aprendizado 1e-4. Na segunda fase (R002), o descongelamento parcial do estágio 4 e do output\_layer entra em ação com taxa de aprendizado 1e-5. Na fase com cabeça auxiliar de relação (R005 em diante), o modelo adiciona uma classificação de onze relações do FIW aplicada apenas aos pares positivos, com peso 0,05 sobre a perda total. A variante R006 introduz a passagem simétrica: cada par é processado nas duas ordens, $(A,B)$ e $(B,A)$, e a perda binária passa a ser a média das duas direções. Essa mudança força consistência aproximada entre $f(A,B)$ e $f(B,A)$, reduzindo atalhos dependentes da ordem do par. As variantes R009 e R010 testam, respectivamente, a remoção dos vetores globais crus da fusão final e a combinação dessa fusão apenas comparativa com taxas de aprendizado diferenciadas para backbone e cabeça. Os resultados aparecem no Capítulo~\ref{cap:resultados}.

\section{Invariância geracional por síntese de idade}
\label{sec:invariancia_geracional}

Comparar um pai de cinquenta anos com um filho de vinte traz um problema conhecido em parentesco facial. O sinal facial fica dominado por mudanças associadas ao envelhecimento, e parte da semelhança de parentesco entre as duas pessoas fica encoberta. Neste trabalho, a estratégia de invariância geracional baseada em síntese de idade foi incorporada apenas ao Modelo 01, que ocupa papel periférico e histórico no estudo. A descrição a seguir registra o procedimento efetivamente associado a esse modelo e delimita sua ausência nos modelos do eixo central de comparação. Cada face passa uma vez por um módulo de síntese de idade que produz três variantes etárias adicionais nas faixas de criança, jovem e idoso, e a comparação entre dois indivíduos agrega quatro pares de imagens em vez de um só.

Dado um par de consulta $(a, b)$, o módulo de síntese gera $a_{\text{cri}}, a_{\text{jov}}, a_{\text{ido}}$ e $b_{\text{cri}}, b_{\text{jov}}, b_{\text{ido}}$. A inferência roda sobre quatro pares diagonais. O par original $(a, b)$ carrega o sinal mais fiel ao indivíduo e recebe o maior peso. Os três pares etários $(a_{\text{cri}}, b_{\text{cri}})$, $(a_{\text{jov}}, b_{\text{jov}})$ e $(a_{\text{ido}}, b_{\text{ido}})$ entram como votos auxiliares que aproximam a comparação de uma mesma faixa etária e reduzem o viés introduzido pela assimetria de idade entre as faces originais. Cada par produz um escore de similaridade $s_k$ pelo mesmo modelo, e o escore final é a combinação convexa

\begin{equation}
s = w_0 \, s_0 + w_{\text{cri}} \, s_{\text{cri}} + w_{\text{jov}} \, s_{\text{jov}} + w_{\text{ido}} \, s_{\text{ido}},
\label{eq:invariancia_geracional}
\end{equation}

com $w_0 + w_{\text{cri}} + w_{\text{jov}} + w_{\text{ido}} = 1$, $w_0 > w_{\text{cri}}, w_{\text{jov}}, w_{\text{ido}}$ e os três pesos secundários iguais entre si. Os valores padrão são $w_0 = 0{,}55$ e $w_k = 0{,}15$ para cada variante etária, e os pesos são reotimizados em validação por grid search no simplex sujeito à restrição $w_0 \ge 0{,}4$.

O módulo de síntese é o SAM \cite{alaluf2022sam}, com pesos pré-treinados externos, descrito no Capítulo~\ref{cap:fundamentacao-teorica}. A síntese roda em pré-processamento, uma única vez por face do banco, e as três variantes ficam em cache em disco. Em inferência o custo extra fica em $4\times$ no número de passagens para frente por par, e em treino o esquema é mantido para que o modelo veja, em cada lote, as quatro variantes do par positivo e do par negativo.

A ativação efetiva do SAM depende da disponibilidade dos pesos pré-treinados de cerca de 4,5 GB e do tempo de pré-processamento das três variantes etárias para o banco completo. No escopo experimental consolidado, apenas o Modelo 01 integra o SAM ao pipeline de treino, e mesmo nele parte das execuções opera sem o módulo ativo. Os Modelos 02, 05, 06 e 12 não usam o módulo de síntese de idade nem o esquema diagonal de quatro pares, e a linha de base B0 também opera só sobre o par original. A invariância geracional, portanto, é tratada como componente do Modelo 01 e não como protocolo uniforme aplicado aos modelos do eixo central de comparação. O Capítulo~\ref{cap:resultados} reporta apenas o regime efetivamente avaliado por cada modelo.

\section{Treinamento}

Todos os modelos supervisionados seguem a mesma receita de treinamento, com pequenas variações por modelo nos casos em que a literatura aponta valores diferentes. A perda padrão é a Supervised Contrastive Loss com distância de cosseno e margem 0,3, escolhida por ter sido a configuração com melhor desempenho no estudo preliminar do Modelo 02. O otimizador é o AdamW, com weight decay de 5e-5 e gradient clipping de norma máxima 1,0.

A taxa de aprendizado tem aquecimento linear de três a oito épocas e decai por cosseno até 1e-7. O tamanho de batch é 16, limite imposto pela memória da GPU usada, e o treinamento usa precisão mista com FP16 quando a plataforma suporta. O dropout é 0,25 e a parada antecipada interrompe o treino após 25 épocas sem melhora na métrica de validação. A semente aleatória é fixada em 42 para amostragem de pares e partições, de modo que cada execução seja reproduzível a partir do comando registrado. A Tabela~\ref{tab:hyperparams} resume os valores padrão.

\begin{table}[h]
\centering
\caption{Receita padrão de hiperparâmetros de treinamento. Valores entre parênteses indicam ajustes finos usados em execuções específicas.}
\label{tab:hyperparams}
\begin{tabular}{ll}
\hline
\textbf{Parâmetro} & \textbf{Valor padrão} \\
\hline
Otimizador & AdamW, weight\_decay = 5e-5 \\
Taxa de aprendizado (pico) & 2e-6 a 5e-6 (M02 R031 em 5e-6, M05 em 3e-4 por LoRA, M12 em 1e-5) \\
Agendador da taxa de aprendizado & Decaimento por cosseno até 1e-7 \\
Aquecimento & 3 a 8 épocas, linear \\
Tamanho de lote & 16 (com acumulação de gradiente onde aplicável) \\
Dropout & 0,25 \\
Recorte de gradiente & max\_norm = 1,0 \\
Precisão mista & AMP FP16 \\
Razão negativos:positivos (treino) & 2:1, amostragem aleatória nas execuções principais \\
Razão negativos:positivos (avaliação) & 1:1, amostragem aleatória fixa \\
Early stopping & paciência 25 épocas \\
Semente & 42 \\
\hline
\end{tabular}
\end{table}

A função de perda varia por modelo e foi escolhida em função da arquitetura e da literatura associada. A Tabela~\ref{tab:perdas_por_modelo} resume a perda principal e as perdas auxiliares de cada modelo no eixo central de comparação.

\begin{table}[h]
\centering
\caption{Função de perda principal e perdas auxiliares por modelo no eixo central de comparação. Pesos auxiliares variam entre execuções e ficam registrados no log experimental.}
\label{tab:perdas_por_modelo}
\begin{tabular}{lll}
\hline
\textbf{Modelo} & \textbf{Perda principal} & \textbf{Perdas auxiliares} \\
\hline
B0  & nenhuma (sem treinamento)                                & --- \\
M02 & Supervised Contrastive (cosseno, margem 0{,}3)            & --- \\
M05 & Supervised Contrastive (cosseno, margem 0{,}3)            & BCE binária; classificação de relação \\
M06 & BCE binária                                               & Supervised Contrastive; classificação de relação \\
M12 & BCE binária; BCE simétrica nas variantes R006+             & classificação de relação nos pares positivos (peso 0,05) \\
\hline
\end{tabular}
\end{table}

\section{Estratégias de plasticidade do backbone}

A flexibilidade do backbone pré-treinado é uma decisão de modelagem importante neste trabalho, e os modelos do eixo central tratam essa decisão de formas distintas. O Modelo 02 é treinado sem congelamento sobre o ViT pré-treinado em ImageNet. O Modelo 05 cobre quatro regimes de plasticidade do DINOv2 que vão de LoRA-only com backbone totalmente congelado até descongelamento completo do encoder, descritos na subseção própria. O Modelo 06 mantém o encoder ViT-B/16 congelado em todas as execuções, e o ganho é absorvido pelo módulo de recuperação e pelos pesos da atenção cruzada sobre os suportes. O Modelo 12 congela os três primeiros estágios do AdaFace IR-101 e libera o último estágio mais a camada de saída, com taxa de aprendizado reduzida para preservar as representações de identidade já consolidadas. A linha de base B0 não é treinada e mantém o backbone integralmente congelado por construção.

O estudo de plasticidade do Modelo 03, que cobre descongelamento completo, congelamento total e duas variantes intermediárias, está descrito no Apêndice~\ref{ap:modelos-perifericos}. O Capítulo~\ref{cap:resultados} usa esse estudo apenas como referência interpretativa, não como parte da tabela principal.

\section{Protocolo de avaliação}

O protocolo de avaliação usa validação cruzada de cinco dobras nos dois bancos, mas com critérios de particionamento diferentes. KinFaceW-I segue o padrão tradicional por pares disjuntos. FIW, por sua vez, é particionado em nível de família, de modo que uma mesma família não aparece simultaneamente em treino, validação e teste dentro de uma dobra. Em ambos os casos, os resultados finais são reportados como média e desvio padrão sobre as cinco dobras, e a escolha de limiar é feita exclusivamente em validação.

Em KinFaceW-I, para cada dobra $i$, o modelo é treinado sobre as outras quatro dobras, validado em uma fração da dobra de treino reservada como conjunto de validação, e testado na dobra $i$. Os hiperparâmetros são fixados antes do início da validação cruzada, e a média sobre as cinco dobras é o resultado reportado, com desvio padrão como medida de variabilidade.

Em FIW, as 761 famílias são distribuídas em cinco dobras família-disjuntas. Em cada iteração, uma dobra funciona como teste, uma segunda dobra como validação e as demais como treino; esse papel gira até que todas as dobras sejam avaliadas como teste. O limiar é escolhido apenas em validação e aplicado uma única vez no teste da dobra correspondente. Quando há múltiplas execuções por modelo, como estudos de congelamento, regime de perda ou amostragem, cada execução aparece com seu identificador no log experimental, e o Capítulo~\ref{cap:resultados} reporta tanto a configuração selecionada para a comparação principal quanto as ablações necessárias para justificar a escolha.

A seleção do limiar de decisão segue o mesmo procedimento nos dois bancos. Após o treino, o modelo produz escores nos dados de validação, e uma busca em grade no intervalo de 0,10 a 0,95 encontra o limiar $\tau^*$ que maximiza F1 na validação. O mesmo $\tau^*$ é aplicado sem reotimização na partição de teste. Essa separação evita o vazamento que ocorreria se o limiar fosse escolhido sobre o próprio conjunto de teste.

A escala do escore varia conforme o modelo, e a busca em grade sobre $[0{,}10;\, 0{,}95]$ só é coerente se todos os escores estiverem no mesmo intervalo aproximado de $[0,1]$. Para garantir isso, cada modelo aplica a transformação adequada antes da busca. O B0 calcula a similaridade cosseno entre embeddings $L^2$-normalizados e mapeia o resultado de $[-1, 1]$ para $[0, 1]$ por $(\cos + 1)/2$. O Modelo 02 e o Modelo 05 produzem logits binários sobre a cabeça final, e o escore usado para a busca é $\sigma(\text{logit})$, a saída do sigmoide. O Modelo 06 e o Modelo 12 também produzem logits binários treinados por BCEWithLogitsLoss, e usam $\sigma(\text{logit})$ pela mesma razão. Assim, em todos os modelos o escore de comparação está em $[0, 1]$, o que torna o limiar escolhido comparável em escala entre métodos.

Em paralelo ao limiar escolhido em validação, três famílias de métrica são reportadas. ROC AUC e Average Precision são independentes de ponto operacional específico, no sentido de que integram desempenho sobre toda a curva de compromisso entre precisão e revocação, e por isso permitem comparação entre métodos sem depender de qual ponto da curva é escolhido. TAR@FAR=0,1, TAR@FAR=0,01 e TAR@FAR=0,001, por outro lado, são calculadas em pontos operacionais fixados pela taxa de falso positivo, e medem a fração de pares positivos corretamente aceitos quando o limiar é definido para produzir falsa aceitação de 10\%, 1\% e 0,1\%, respectivamente. Essas métricas correspondem a regimes de alta precisão e são relevantes em aplicações forenses, em que falsos positivos têm custo alto. Acurácia, F1, precisão e revocação, quando apresentadas, ficam dependentes do limiar $\tau^*$ escolhido em validação e por isso são interpretadas em conjunto com o ponto operacional declarado.

A linha de base B0 segue a mesma regra de seleção de limiar em validação. Os modelos exploratórios M05 e M06 usam o regime declarado em suas respectivas subseções, em que o M05 cobre quatro regimes de plasticidade do backbone descritos antes (LoRA-only, descongelamento parcial, descongelamento completo e híbrido), e o M06 mantém o encoder ViT-B/16 congelado em sua execução principal. A galeria do Modelo 06 é reconstruída a cada dobra a partir da partição de treino correspondente, sem qualquer informação da validação ou do teste.

\section{Métricas}

O protocolo contempla três grupos de métricas. As métricas binárias avaliam a tarefa de verificação propriamente dita e cobrem acurácia, precisão, revocação, F1, ROC AUC, Average Precision e taxas de aceitação verdadeira em três taxas fixas de falso positivo, TAR@FAR=0,1, TAR@FAR=0,01 e TAR@FAR=0,001.

O ROC AUC é a métrica principal por ser independente do limiar escolhido. As métricas TAR@FAR são relevantes porque medem o desempenho em regime de alta precisão, mais próximo do uso prático em aplicações forenses, onde falsos positivos têm custo alto.

Por relação, reportam-se acurácia e F1 condicionais ao tipo de parentesco do par. Esse recorte é importante porque a literatura mostra desempenho muito heterogêneo entre relações, com pais e filhos atingindo valores altos e avós e netos historicamente mais difíceis.

Para as linhas de base com VLMs, descritas na próxima seção e que avaliam classificação em onze classes em vez de verificação binária, aplicam-se métricas multiclasse específicas. Para essa tarefa são reportadas acurácia, macro-precision, macro-recall, macro-F1 e matriz de confusão.

\section{Linhas de base com modelos multimodais}

As linhas de base com VLMs entram como medida do que um modelo multimodal generalista consegue inferir sobre parentesco facial sem treinamento específico no domínio. O Codex é avaliado em três regimes progressivos, e o Claude Sonnet é avaliado em um recorte piloto zero-shot.

\subsection{Configuração comum}

A tarefa é a classificação multiclasse das onze relações positivas do FIW. A entrada é uma imagem composta lado a lado, formada por duas faces do par a ser classificado, e a saída válida é uma das onze relações. As amostras são extraídas da partição de teste do FIW com seleção estratificada quase uniforme entre as classes positivas. Os modelos não recebem qualquer treinamento sobre FIW e operam apenas a partir do prompt textual e da imagem composta. Respostas fora do conjunto de onze relações, inclusive respostas de não-parentesco, são contabilizadas como erro de classificação. Por usarem tarefa diferente dos modelos supervisionados, os números reportados para VLMs são comparados qualitativamente, não métrica contra métrica.

\subsection{Codex zero-shot}

A primeira linha de base é o gpt-5.4-mini, executado via Codex CLI com \texttt{reasoning\_effort=low}. A amostra é de 750 pares positivos balanceados entre as onze classes, totalizando 1.500 imagens. O prompt é direto, descreve as onze classes possíveis, e pede a classificação da relação entre as duas faces.

\subsection{Claude Sonnet zero-shot}

A segunda linha de base é o claude-sonnet-4-6, avaliado por visão direta sobre a mesma estrutura de prompt do Codex. A avaliação do Claude usa um recorte menor de 75 pares balanceados, retirados do mesmo conjunto de teste. Isso permite comparação direta com o Codex no recorte de 75 pares, e comparação qualitativa sobre os 750 pares.

\subsection{Adaptação ao domínio por prompt no Codex}

Para modificar o comportamento do VLM sem fine-tuning, este regime combina três mecanismos. O few-shot in-context learning entra com sete a onze exemplos retirados da partição de validação do FIW, escolhidos para cobrir as classes mais difíceis. Em torno desses exemplos, um prompt estruturado obriga o modelo a estimar, antes da resposta final, a diferença geracional aparente, o gênero de cada face e qual face parece mais velha. Uma calibração leve treinada em 110 pares de validação ajusta a resposta final a partir dessas saídas auxiliares.

Quatro configurações são comparadas em validação, e a melhor é avaliada nos 750 pares de teste, os mesmos do zero-shot, o que permite comparação direta entre os dois regimes. A adaptação aqui é estritamente inferencial e não envolve atualização de pesos do modelo.

\subsection{Prompts direcionados por treinamento}

Um terceiro regime de VLMs é considerado, baseado em curadoria iterativa de exemplos e estruturas de prompt otimizadas sobre um subconjunto de validação. A diferença em relação à adaptação por prompt é o ciclo de iteração, em que o conjunto de exemplos few-shot e a estrutura do prompt são ajustados em rodadas sucessivas a partir do desempenho observado em validação. Esse regime não foi executado sobre os 750 pares do teste completo. Ele aparece apenas como recorte diagnóstico de 12 pares das quatro relações de avós e netos, usado para investigar se prompts explicitamente orientados por diferença geracional corrigiriam o principal erro do VLM.

\begin{table}[h]
\centering
\caption{Especificação das linhas de base com modelos multimodais.}
\label{tab:vlms_metodologia}
\begin{tabular}{lll}
\hline
\textbf{Modelo} & \textbf{Regime} & \textbf{Amostra de teste} \\
\hline
Codex (gpt-5.4-mini) & Zero-shot                       & 750 pares balanceados \\
Codex (gpt-5.4-mini) & Adaptação por prompt (7-shot)    & 750 pares (mesmos do zero-shot) \\
Codex (gpt-5.4-mini) & Prompts direcionados por treino  & 12 pares de avós e netos (diagnóstico) \\
Claude Sonnet 4.6   & Zero-shot direto                  & 75 pares (piloto) \\
\hline
\end{tabular}
\end{table}

\section{Hardware e software}

Todos os experimentos foram executados em uma única máquina com GPU AMD Radeon RX 6750 XT, arquitetura RDNA2 com identificador gfx1031, 12 GB de memória GDDR6. A plataforma de software é o ROCm 5.7 com PyTorch 2.3.1 compilado para essa versão do ROCm. As variáveis de ambiente \texttt{HSA\_OVERRIDE\_GFX\_VERSION=10.3.0}, \texttt{MIOPEN\_FIND\_MODE=FAST} e \texttt{HSA\_FORCE\_FINE\_GRAIN\_PCIE=1} são definidas antes da inicialização do framework para garantir compatibilidade e desempenho razoável. O treinamento usa precisão mista, FP16 nas operações suportadas e FP32 nas demais.

Este hardware impôs restrições reais nas escolhas metodológicas. O tamanho máximo de lote com backbones descongelados é 16, e qualquer treinamento que demande mais de cerca de 10 GB de memória excede o limite disponível. O Modelo 03, com aproximadamente 176 milhões de parâmetros, inviabiliza descongelamento completo em lote maior nesta GPU. O Modelo 05 viabiliza-se nesta configuração porque o backbone permanece congelado e os adaptadores LoRA mantêm a memória ativa em níveis administráveis. O Modelo 06 mantém a galeria de pares positivos principalmente em RAM da máquina, fora da GPU, para liberar VRAM ao encoder e à atenção cruzada de recuperação.

\section{Reprodutibilidade}

Cada modelo tem seu próprio diretório no repositório, com pipeline de execução em ROCm e em CUDA. Um script orquestrador, \texttt{AMD/run\_pipeline.sh}, encadeia treino, teste e avaliação, gera saídas em diretórios numerados automaticamente (\texttt{output/001/}, \texttt{output/002/} e assim por diante, sem sobrescrita) e produz, ao final de cada execução, um arquivo \texttt{RUN\_LOG.md} no diretório do modelo. O \texttt{RUN\_LOG.md} registra o comando exato com todas as variáveis de ambiente, a configuração completa em formato tabular, a trajetória de treinamento época a época, e as métricas finais de teste por relação. Qualquer execução reportada no Capítulo 5 pode ser recriada a partir desse registro.

As métricas em formato JSON, as figuras de ROC e matriz de confusão, e os pontos salvos de melhor desempenho ficam em \texttt{output/<id>/results/} e \texttt{output/<id>/checkpoints/} de cada modelo. As sementes para amostragem de pares ficam fixadas em 42; as dobras família-disjuntas do FIW são geradas deterministicamente a partir do Track-I do RFIW, e as cinco dobras de KinFaceW-I são geradas deterministicamente e identicamente entre execuções. Os artefatos dos VLMs ficam em \texttt{data/codex\_vlm\_*} e \texttt{data/claude\_vlm\_*}, incluindo prompt, configuração, manifesto de pares, respostas brutas, métricas agregadas e matriz de confusão.
